\documentclass[11pt,a4paper]{article}
\usepackage[utf8]{inputenc}
\usepackage[T1]{fontenc}
\usepackage[brazil]{babel}
\usepackage{lmodern}
\usepackage{microtype}
\usepackage{geometry}
\usepackage{tabularx}
\usepackage{booktabs}
\usepackage{array}
\usepackage{enumitem}
\usepackage{fancyhdr}
\usepackage{titlesec}
\usepackage{hyperref}
\usepackage{xurl}
\usepackage{xcolor}
\usepackage{listings}
\geometry{left=2.25cm,right=2.25cm,top=2.25cm,bottom=2.1cm}
\definecolor{marinho}{HTML}{193D62}
\definecolor{cinza}{HTML}{65717D}
\definecolor{fundo}{HTML}{F2F5F9}
\definecolor{laranja}{HTML}{C46B34}
\hypersetup{colorlinks=true,urlcolor=marinho,linkcolor=marinho,pdftitle={Atividade investigativa II - Comparacao de deadlocks entre Linux Windows macOS},pdfauthor={Sistemas Operacionais}}
\titleformat{\section}{\large\bfseries\color{marinho}}{\thesection.}{.6em}{}
\titleformat{\subsection}{\normalsize\bfseries\color{marinho}}{\thesubsection.}{.6em}{}
\titlespacing*{\section}{0pt}{10pt}{4pt}
\titlespacing*{\subsection}{0pt}{7pt}{3pt}
\pagestyle{fancy}
\fancyhf{}
\fancyhead[L]{\footnotesize\color{cinza}Sistemas Operacionais | Atividade investigativa II}
\fancyhead[R]{\footnotesize\color{cinza}Linux \;\textbullet\; Windows \;\textbullet\; macOS}
\fancyfoot[C]{\footnotesize\thepage}
\renewcommand{\headrulewidth}{.3pt}
\setlength{\headheight}{15pt}
\setlength{\parindent}{0pt}
\setlength{\parskip}{4pt}
\setlist[enumerate]{leftmargin=*,itemsep=2.5pt,topsep=2pt}
\setlist[itemize]{leftmargin=*,itemsep=2pt,topsep=2pt}
\lstdefinestyle{terminal}{basicstyle=\ttfamily\small,backgroundcolor=\color{fundo},frame=single,rulecolor=\color{fundo},xleftmargin=5pt,xrightmargin=5pt,aboveskip=4pt,belowskip=4pt,showstringspaces=false,breaklines=true}
\newcommand{\emphbox}[1]{\vspace{2pt}\noindent\colorbox{fundo}{\parbox{\dimexpr\linewidth-2\fboxsep}{#1}}\vspace{2pt}}
\begin{document}
\begin{center}
{\Large\bfseries\color{marinho} Atividade Investigativa II}\\[3pt]
{\large\bfseries Deadlocks: uma comparação entre sistemas operacionais}\\[5pt]
{\normalsize Linux, Windows e macOS \;|\; APIs de bloqueio de arquivos}\\[5pt]
{\footnotesize Ciência da Computação \;--\; Sistemas Operacionais}
\end{center}
\vspace{1mm}
{\color{marinho}\hrule height .9pt}
\vspace{3mm}
\textbf{Integrante(s):} \rule{11.4cm}{.4pt}\\[6pt]
\textbf{Turma:} \rule{5cm}{.4pt}\hfill\textbf{Data:} \rule{4cm}{.4pt}

\section{Proposta da atividade}
Esta é uma \textbf{nova atividade, independente da atividade anterior sobre deadlocks no Linux}. O objetivo é comparar \textbf{como três sistemas operacionais se comportam diante de pedidos de bloqueio de arquivos em ordens opostas}, utilizando mecanismos próprios de cada plataforma.

O trabalho exige: \textbf{formular hipóteses, ler os códigos, executar os testes, registrar evidências e interpretar os resultados}. Não assumam que chamadas parecidas tenham a mesma semântica, nem que os sistemas operacionais necessariamente produzam o mesmo resultado.

\emphbox{\textbf{Regra da investigação:} este roteiro explica o que as APIs solicitam, mas \textbf{não antecipa o desfecho dos testes}. Diferenciem sempre o comportamento observado de suas hipóteses e de informações obtidas na documentação.}

\section{O que é igual e o que é diferente?}
Em todos os códigos, duas unidades de execução disputam dois arquivos, \texttt{A.txt} e \texttt{B.txt}: uma solicita A e depois B; a outra solicita B e depois A. Há um intervalo de dez segundos entre as solicitações.

\renewcommand{\arraystretch}{1.18}
\begin{tabularx}{\textwidth}{>{\bfseries}p{2.55cm}p{2.85cm}X}
\toprule
Sistema / API & Unidades de execução & O que a chamada solicita\\
\midrule
Linux: \texttt{flock()} & Duas \textit{threads} & Bloqueio consultivo exclusivo do arquivo; associado à descrição de arquivo aberto no Linux.\\
Linux: \texttt{fcntl()} & Dois processos & Bloqueio tradicional de escrita sobre uma região; associado ao processo no modelo POSIX.\\
macOS: \texttt{flock()} & Duas \textit{threads} & Bloqueio exclusivo solicitado sobre o arquivo por meio de uma API POSIX/BSD.\\
macOS: \texttt{fcntl()} & Dois processos & Bloqueio de escrita sobre uma região por \texttt{F\_SETLKW}, com espera caso necessário.\\
Windows:\newline \texttt{LockFileEx()} & Duas \textit{threads} & Bloqueio exclusivo de um intervalo de bytes por meio de um \textit{handle} Win32.\\
\bottomrule
\end{tabularx}

\textbf{Observação conceitual:} abrir com \texttt{O\_RDWR} no POSIX ou solicitar acesso de leitura/escrita com \texttt{CreateFileA()} no Windows não é o mesmo que adquirir os bloqueios solicitados posteriormente. As regras de compartilhamento de abertura do Windows também não devem ser confundidas com \texttt{LockFileEx()}.

\newpage
\section{Antes de executar: análise dos códigos}
Encontrem, em \textbf{cada programa}, as chamadas responsáveis pela abertura do primeiro arquivo, pela aquisição do primeiro bloqueio, pela pausa, pela abertura do segundo arquivo, pela solicitação do segundo bloqueio e pela liberação de recursos.

\begin{enumerate}[label=H\arabic*.]
\item Para cada API, qual resultado vocês esperam quando as duas unidades solicitarem o segundo bloqueio? Registrem a justificativa \textbf{antes} da execução.
\item A troca de \textit{threads} por processos pode alterar o significado dos bloqueios? Expliquem sua hipótese.
\item Quais diferenças esperam encontrar entre Linux e macOS usando os nomes \texttt{flock()} e \texttt{fcntl()}? Nomes iguais garantem implementação idêntica?
\item No Windows, o fato de bloquear apenas \textbf{um byte} pode ser relevante para o conflito entre as threads? Justifiquem.
\end{enumerate}

\section{Experimentos no Linux}
\textbf{Diretório:} \texttt{Atividade\_Comparativa/Linux/}. Utilizem um terminal Linux e um sistema de arquivos local.
\begin{lstlisting}[style=terminal]
make
make teste_flock
make teste_fcntl
\end{lstlisting}
\textbf{Programa 1 -- \texttt{01\_flock\_threads.c}:} localizem a criação e a espera das threads, pelas chamadas \texttt{pthread\_create()} e \texttt{pthread\_join()}. Identifiquem também \texttt{open()}, \texttt{flock()} e \texttt{sleep(10)}.

\textbf{Programa 2 -- \texttt{02\_fcntl\_processos.c}:} verifiquem como o Makefile inicia dois processos em ordens opostas. Examinem \texttt{struct flock}, \texttt{F\_WRLCK}, \texttt{F\_SETLKW} e os descritores abertos por cada processo.

\emphbox{\textbf{Atenção:} no Linux, o comando \texttt{timeout} do Makefile limita externamente o tempo de observação. \textbf{Não o interpretem como uma detecção ou recuperação feita pelo kernel}. Verifiquem se restaram processos em execução antes de repetir o teste.}

\section{Experimentos no macOS}
\textbf{Diretório:} \texttt{Atividade\_Comparativa/MacOS/}. Utilizem o Terminal do macOS com as ferramentas de linha de comando do Xcode instaladas.
\begin{lstlisting}[style=terminal]
make
./01_flock_threads
./02_fcntl_processos
\end{lstlisting}
\textbf{Programa 3 -- \texttt{01\_flock\_threads.c}:} localizem as duas chamadas a \texttt{pthread\_create()}, as aberturas separadas e os pedidos de bloqueio exclusivo.

\textbf{Programa 4 -- \texttt{02\_fcntl\_processos.c}:} localizem \texttt{fork()}, \texttt{getpid()}, \texttt{struct flock} e a chamada a \texttt{F\_SETLKW}. Observem por que as aberturas ocorrem nos processos filhos.

Executem \textbf{um teste de cada vez}. Caso um teste não termine, usem \texttt{Ctrl+C} e registrem que a interrupção foi manual, não uma conclusão espontânea do programa.

\newpage
\section{Experimento no Windows}
\textbf{Diretório:} \texttt{Atividade\_Comparativa/Windows/}. Este é um experimento \textbf{Win32 nativo}; executá-lo dentro do WSL testaria outro ambiente.

\textbf{Com MinGW-w64 (CMD ou terminal apropriado):}
\begin{lstlisting}[style=terminal]
compilar_mingw.bat
01_lockfileex_threads.exe
\end{lstlisting}
\textbf{Alternativa com MSVC:} no \textit{Developer Command Prompt}, executem:
\begin{lstlisting}[style=terminal]
compilar_msvc.bat
01_lockfileex_threads.exe
\end{lstlisting}

\textbf{Programa 5 -- \texttt{01\_lockfileex\_threads.c}:} identifiquem:
\begin{itemize}
\item a abertura: \texttt{CreateFileA()};
\item as flags de compartilhamento: \texttt{FILE\_SHARE\_READ} e \texttt{FILE\_SHARE\_WRITE};
\item os bloqueios e desbloqueios: \texttt{LockFileEx()} e \texttt{UnlockFileEx()};
\item a espera e o tratamento de erros: \texttt{Sleep()} e \texttt{GetLastError()};
\item a espera pelas duas threads: \texttt{WaitForMultipleObjects()}.
\end{itemize}

\emphbox{\textbf{Detalhe para investigar:} o código Windows solicita um bloqueio exclusivo de \textbf{um byte a partir do deslocamento zero}, usando um \texttt{HANDLE} específico. Os exemplos de \texttt{fcntl()} usam uma região que se estende até o fim do arquivo. Observem o que é comparável e o que não é.}

\section{Registro experimental}
Antes de discutir o resultado, anotem para \textbf{cada experimento}: sistema operacional e versão, sistema de arquivos local quando conhecido, compilador, comandos executados, mensagens relevantes e como a execução terminou.

\begin{tabularx}{\textwidth}{>{\bfseries}p{2.8cm}X X}
\toprule
Teste & Evidências observadas & Finalização / erro (se houver)\\
\midrule
Linux: \texttt{flock()} & \rule{0pt}{1.0cm} & \\
Linux: \texttt{fcntl()} & \rule{0pt}{1.0cm} & \\
macOS: \texttt{flock()} & \rule{0pt}{1.0cm} & \\
macOS: \texttt{fcntl()} & \rule{0pt}{1.0cm} & \\
Windows: \texttt{LockFileEx()} & \rule{0pt}{1.0cm} & \\
\bottomrule
\end{tabularx}

Se alguma chamada retornar erro, registrem o \textbf{código numérico ou simbólico} e pesquisem seu significado no manual oficial da plataforma. Não atribuam automaticamente todo erro a um deadlock.

\section{Condições de Coffman}
Para cada estado de espera observado, identifiquem: \textbf{exclusão mútua}, \textbf{posse e espera}, \textbf{ausência de preempção} e \textbf{espera circular}. Quando fizer sentido, representem o estado por um grafo de espera. Distingam \textbf{um estado de espera circular}, \textbf{a detecção de deadlock} e \textbf{a recuperação de um impasse}.

\section{Questões comparativas}
Respondam com evidências dos códigos e das execuções. Ao utilizar documentação, identifiquem a fonte consultada.
\begin{enumerate}[label=Q\arabic*.]
\item Quais diferenças existem entre \texttt{open(..., O\_RDWR)}, \texttt{CreateFileA()} e as operações efetivas de bloqueio utilizadas nos programas?
\item Que entidade cada API considera para associar o bloqueio: processo, descrição de arquivo aberta ou \textit{handle}? Que partes dos códigos ajudam a verificar isso?
\item No Linux, que diferenças vocês observaram entre \texttt{flock()} e \texttt{fcntl()}? Como os códigos tratam possíveis erros e a liberação dos recursos?
\item Comparando Linux e macOS, as mesmas APIs produziram o mesmo comportamento? Quais evidências sustentam a resposta?
\item Como o comportamento do Windows se compara ao das duas plataformas POSIX? Que cuidados são necessários para comparar bloqueios de arquivo inteiro com bloqueios de intervalos de bytes?
\item Houve casos de espera, falha, interrupção manual ou término espontâneo? Expliquem \textbf{quem} produziu cada evento: o programa, o kernel, a API ou um comando externo.
\item Em quais experimentos foi possível identificar as condições de Coffman? Um erro retornado por uma chamada implica, por si só, recuperação do estado da aplicação?
\item Proponham uma alteração \textbf{pequena} no código para evitar a aquisição em ordens inversas. Justifiquem sua proposta e, se puderem, realizem um teste adicional em cópia do programa.
\end{enumerate}

\section{O que entregar}
Entreguem \textbf{um novo relatório em PDF}, separado da atividade original do Linux, contendo: identificação; hipóteses H1--H4; tabela de observações preenchida; saídas representativas; respostas Q1--Q8; análise das condições de Coffman; e uma conclusão comparativa.

\textbf{Acesso a sistemas diferentes:} executem os testes nos ambientes aos quais realmente tenham acesso. Para uma plataforma indisponível, façam pesquisa documental e deixem explícito que \textbf{não houve execução nativa}. Nunca apresentem como observado um resultado apenas pesquisado ou presumido.

\section{Fontes para consulta}
\begin{itemize}
\item \textbf{Linux:} \url{https://man7.org/linux/man-pages/man2/flock.2.html} e \url{https://man7.org/linux/man-pages/man2/fcntl_locking.2.html}.
\item \textbf{macOS:} manuais \texttt{man 2 flock} e \texttt{man 2 fcntl}; \url{https://developer.apple.com/library/archive/documentation/System/Conceptual/ManPages_iPhoneOS/man2/fcntl.2.html}.
\item \textbf{Windows:} \url{https://learn.microsoft.com/en-us/windows/win32/api/fileapi/nf-fileapi-lockfileex} e documentação de \texttt{CreateFileA()} e \texttt{GetLastError()}.
\end{itemize}
\end{document}
